% Work-only exact partial-summand theorem for independent audit.
\section*{Two paired handle tangencies in genus three}

Let $V$ be symplectic of dimension six with basis
$a,b,c,d,e,f$, paired as $(a,b),(c,d),(e,f)$.
Let $U_{3,6}\subset\mathfrak h_{3,6}$ and
$q_6:U_{3,6}\to Q_6$ be the actual universal
arithmetic image and quotient of F.3.24.
Write $P_{\mathrm{all}}$ for the subspace of
special degree-six derivations tangent to the
first four basis lines:
\[
 D(x)\in[\mathbb L_6(V),x]\qquad(x=a,b,c,d).
\]
The values on $e,f$ are unrestricted except by
the specialness equation. Then
\[
                         q_6(U_{3,6}\cap P_{\mathrm{all}})=0.
                                                               \tag{PT.1}
\]
This gives a partial-summand statement on the original
six-dimensional space. No arithmetic curve quotient
or compatibility of $q_6$ with a restriction map is used.

\paragraph{1. The actual geometric quotient and all invariant functionals.}
The original exact sequence (4.7.1) of Nakamura,
printed p.~117~\cite{AuditNakamura1996}, identifies the universal quotient
as the quotient of the universal image by its
geometric image. The last paragraph of that page
also explicitly places the topological Johnson
image in the geometric image. Under the strict
punctured-action map of F.3.24, this gives
\[
                    J_6\subset U_{3,6},\qquad q_6(J_6)=0.
                                                               \tag{PT.2}
\]
Here $J_6$ denotes the actual degree-six geometric
Johnson image in the canonical special derivations.
In particular every sixfold Lie bracket of
degree-one Johnson derivations belongs to this
specified kernel. These are actual geometric
image vectors, not Lie polynomials merely
declared to be arithmetic.

Choose a nonzero basis of the one-dimensional $Q_6$,
whose dimension was verified in F.3.27, to regard $q_6$
as a scalar-valued map. By complete reducibility, the
symplectic-invariant map $q_6$ extends to a symplectic-invariant linear
functional $F$ on $\mathfrak h_{3,6}$, retaining
its zero value on $J_6$. The natural tensor map
\[
 \mathfrak h_{3,6}\hookrightarrow V\otimes\mathbb L_7(V)
                      \hookrightarrow V^{\otimes8}
\]
sends $D$ to
\[
 a\otimes D(b)-b\otimes D(a)
 +c\otimes D(d)-d\otimes D(c)
 +e\otimes D(f)-f\otimes D(e).
\]
Extend $F$ once more, equivariantly, to $V^{\otimes8}$.
Every invariant functional there is a linear
combination of the 105 complete symplectic
contractions indexed by perfect matchings of
the eight tensor slots. This is precisely the
all-genus surjectivity of Morita--Sakasai--Suzuki,
original v2, Proposition~2.1, p.~8, with the
dual contraction formula on p.~7~\cite{AuditMoritaSakasaiSuzuki2015v2}. Its bijectivity
requires $g\geq4$ here, but its stated surjectivity
holds for every $g$, including $g=3$.
The program uses the unsigned products of the
four pairings; these differ from that paper's
convention by an invertible sign on each of
the 105 coordinates, which changes neither
their span nor the image inclusion below.

Let $K:\mathfrak h_{3,6}\to\Q_\ell^{105}$ be
the vector of those contractions. To show
$F(D)=0$ it therefore suffices to show
\[
                    K(D)\in K(J_6).                       \tag{PT.3}
\]
This implication uses a functional on the original
derivation space. It makes no identification with
an arithmetic quotient belonging to a smaller curve.

\paragraph{2. The complete balanced partial-tangent space.}
Give the full diagonal similitude torus the weights
$a:\alpha_1,b:t/\alpha_1,c:\alpha_2,d:t/\alpha_2,
e:\alpha_3,f:t/\alpha_3$, and let $P$ be the
$t^3$ part of $P_{\mathrm{all}}$.
The specialness equation has balanced output
multidegrees $(n_1,n_1,n_2,n_2,n_3,n_3)$ with
$n_1+n_2+n_3=4$.

For a constrained generator $x\in\{a,b,c,d\}$,
the column space consists exactly of
$[A_x,x]$, with $A_x$ in the corresponding
degree-six Lyndon basis. For an unconstrained
generator $x\in\{e,f\}$, the column space
is the complete degree-seven Lyndon basis
of the required multidegree. Applying
$D\omega$ gives the following full block table:
\[
\begin{array}{c|r|r|r|r}
n&\text{rows}&\text{columns}&\text{rank}&\text{nullity}\\\hline
(4,0,0),(0,4,0)&8&6&6&0\\
(0,0,4)&8&10&8&2\\
(3,1,0),(1,3,0)&140&66&66&0\\
(3,0,1),(0,3,1)&140&100&99&1\\
(1,0,3),(0,1,3)&140&126&123&3\\
(2,2,0)&312&120&119&1\\
(2,0,2),(0,2,2)&312&240&236&4\\
(2,1,1),(1,2,1)&1260&660&657&3\\
(1,1,2)&1260&840&834&6.
\end{array}                                               \tag{PT.4}
\]
For rows listing two multidegrees, the displayed
sizes and nullity apply to each block. They sum
to $\dim P=31$.

These are exact rational ranks. The complete
program \texttt{check\_genus\_three\_two\_pair\_partial.py}
performs elimination modulo~101 for the lower
bounds and lifts each kernel vector to an
integer vector, checking every original
specialness equation over the integers.
The nonzero multiplier of each lift modulo~101
preserves the independence of the modular
free-coordinate kernel basis. The upper and
lower bounds therefore agree over $\Q$.
This argument does not identify a modular
kernel with a rational one without checking
the lifted equations.

For all twenty degree-one derivations
$\delta_{uvw}$ of F.3.24, form
$\operatorname{ES}_7([D,\delta_{uvw}])$.
The exact matrix on the 31 basis vectors has
3920 nonzero coordinate rows and rank~15.
The program supplies sixteen independent
integer kernel vectors and verifies all their
entries against the original integer matrix.
Thus the complete subspace
\[
 P^0=\{D\in P:\operatorname{ES}_7([D,\delta])=0
                      \text{ for every }\delta\in\Lambda^3V\}
\]
has dimension~16 over $\Q$ and over $\Q_\ell$.
The actual geometric-normalizer equation of
F.3.24 places $U_{3,6}\cap P$ inside $P^0$.

\paragraph{3. Exact contraction image in two specified geometric brackets.}
Write brackets right-nested, so that
$[D_1,D_2,\ldots,D_6]$ means
$[D_1,[D_2,[D_3,[D_4,[D_5,D_6]]]]]$.
Define the actual geometric degree-six elements
\[
\begin{split}
 Z_1&=[\delta_{cdf},\delta_{cef},\delta_{aef},
          \delta_{def},\delta_{bef},\delta_{cde}],\\
 Z_2&=[\delta_{bce},\delta_{adf},\delta_{adf},
          \delta_{abd},\delta_{bce},\delta_{abc}].
\end{split}                                               \tag{PT.5}
\]
They lie in $J_6$ by its graded Lie bracket.
The contraction vectors $K(Z_1),K(Z_2)$ are
independent: on the matchings
$(12)(34)(56)(78)$ and $(12)(36)(47)(58)$
their two-by-two determinant is $-10752$.

For each of the sixteen exact basis vectors
$D_j$ of $P^0$, the program computes its
105 integer contractions and solves for the
two coefficients using that nonzero minor.
It then verifies the resulting identity in
every one of the 105 coordinates over the
integers after multiplying by $-10752$.
The complete list of integer coefficient
triples is retained in the output. Hence
\[
       K(P^0)\subset\operatorname{span}_{\Q}
                         \{K(Z_1),K(Z_2)\}\subset K(J_6).
                                                               \tag{PT.6}
\]
The nonzero rational determinant remains invertible
in every $\Q_\ell$; no integral division is needed.
Both the source kernel and the image inclusion
are exact; a finite-field rank coincidence
alone is not used for (PT.6).

For $D\in U_{3,6}\cap P$, choose the invariant
extension $F$ above. Express it as a linear
combination of the 105 contractions. Equation
(PT.6) then gives
$F(D)=\alpha F(Z_1)+\beta F(Z_2)=0$
by (PT.2). Since $F(D)=q_6(D)$ on the actual
universal submodule, the balanced assertion follows.

Finally $P_{\mathrm{all}}$ is stable under the
full diagonal torus despite imposing no
constraint on $e,f$. The explicit weight
projector of F.3.25 preserves its intersection
with $U_{3,6}$, projects to $P$, and preserves
$q_6$. This proves (PT.1) for all degrees of
the unconstrained torus weights as well.

\paragraph{Exact scope and the one-pair test.}
The theorem imposes four actual line tangencies
spanning a nondegenerate four-dimensional
symplectic summand of this genus-three homology.
It retains both complementary generators and
all mixed Johnson equations. It makes no
assertion that arbitrary CM quotient lines
necessarily supply that nondegenerate summand.
A separate discovery calculation with only
one paired handle has an additional contraction
direction modulo~101; it has not supplied an
arithmetic class or a proof of a converse.
That calculation is preserved as work provenance
and is not used in the theorem.
